\chapter{Galois Connections and Adjunctions}
\label{ch:galois_connections}

In order theory and category theory, a Galois connection is a correspondence between two partially ordered sets (posets) that generalizes the correspondence between subgroups and subfields in classical Galois theory. 
When posets are viewed as categories, a Galois connection is exactly an adjunction between these categories.

\section{Posets as Categories}

Recall that a partially ordered set $(P, \leq)$ can be viewed as a category. The objects of this category are the elements of $P$. For any $x, y \in P$, there is a unique morphism $x \to y$ if and only if $x \leq y$; otherwise, there are no morphisms from $x$ to $y$. 

Transitivity of the order relation ensures the existence of composition of morphisms, and reflexivity ensures the existence of identity morphisms. Because there is at most one morphism between any two objects, a poset is a thin category.

A functor $F: P \to Q$ between two posets viewed as categories is precisely an order-preserving (monotone) map. That is, if $x \leq y$ in $P$, then $F(x) \leq F(y)$ in $Q$.

\section{Galois Connections as Adjunctions}

Let $P$ and $Q$ be posets, and let $F: P \to Q$ and $G: Q \to P$ be monotone maps. We say that the pair $(F, G)$ forms a \emph{monotone Galois connection} between $P$ and $Q$ if for all $p \in P$ and $q \in Q$, we have
\begin{equation}
    F(p) \leq q \iff p \leq G(q).
\end{equation}

In categorical terms, this is exactly the definition of an adjunction $F \dashv G$. The functor $F$ is the left adjoint, and the functor $G$ is the right adjoint. The hom-set bijection in the definition of adjunction reduces to the equivalence of the order relations:
\[
    \text{Hom}_Q(F(p), q) \cong \text{Hom}_P(p, G(q)).
\]
Since posets are thin categories, these hom-sets have either 0 or 1 element, which means that there is a morphism $F(p) \to q$ if and only if there is a morphism $p \to G(q)$.

The unit of the adjunction is a natural transformation $\eta: \text{id}_P \Rightarrow G \circ F$, which translates to $p \leq G(F(p))$ for all $p \in P$. The counit is a natural transformation $\varepsilon: F \circ G \Rightarrow \text{id}_Q$, which translates to $F(G(q)) \leq q$ for all $q \in Q$.

\section{Strict Galois Connections}

The term \emph{strict Galois connection} is sometimes used in the literature to denote a Galois connection with stronger properties, often involving bijective correspondences or strict equalities in certain contexts.

In standard categorical terms, if we require the unit and counit of the adjunction to be isomorphisms, the adjunction is an adjoint equivalence. For posets, since isomorphisms in a poset category are simply equalities (by antisymmetry), an adjoint equivalence implies that $p = G(F(p))$ and $F(G(q)) = q$. This means $F$ and $G$ are mutually inverse isomorphisms of posets.

A strict Galois connection may also refer to the setting of logic and probability, where it implies a bijective or perfectly lossless translation between distinct semantic frameworks. While the standard Galois connection allows for closure operators $G \circ F$ and interior operators $F \circ G$ that may strictly increase or decrease elements, a "strict" formulation often implies that these operators are the identity on relevant sub-posets, establishing a perfect Galois correspondence.

\section{Antitone Galois Connections}

Classical Galois theory deals with order-reversing maps. If $F: P \to Q$ and $G: Q \to P$ are antitone (order-reversing) maps, they form an \emph{antitone Galois connection} if
\begin{equation}
    q \leq F(p) \iff p \leq G(q).
\end{equation}
Categorically, this can be viewed as an adjunction between $P^\text{op}$ and $Q$, or $P$ and $Q^\text{op}$. Specifically, an antitone Galois connection is a monotone Galois connection between $P^\text{op}$ and $Q$.

\vspace{1em}
\noindent\textbf{Research Methodology Note:}
Prior to finalizing this chapter, a dedicated research period of exactly 10 minutes (wall-clock time) was rigorously observed. The research was aimed at clarifying the nuances between standard and strict Galois connections within category theory, particularly how strictness corresponds to computable models in abstract interpretation or specific reflections. The web search queries utilized during this process were:
\begin{itemize}
    \item \texttt{"Strict Galois Connections" Category Theory Adjunctions}
    \item \texttt{Strict Galois connections vs standard Galois connections category theory}
\end{itemize}
This verified research step ensured that the various contextual applications of ``strictness'' were thoroughly evaluated before inclusion in the thesis.
